\section{Introduction}\label{sec:intro}

A directed cycle cover assigns each vertex a successor and visits every
vertex exactly once, possibly in several disjoint cycles. Give each
directed edge a variable. The weight of a cover is the product of the
variables on its edges. The permanent sums the weights of all covers. The Hamiltonian
cycle polynomial sums those that visit all vertices in a single tour.
Can an arithmetic circuit evaluate either sum with polynomially many
operations? Counting the terms gives little indication of the answer.
The determinant assigns a sign to each cover and has a polynomial-size
circuit. A circuit can reuse intermediate results and cancel terms in
ways that the expansion never reveals.

We prove a superpolynomial circuit lower bound for Hamiltonian cycle
over every field and for the permanent over every field of
characteristic different from two. The exception for the permanent is
necessary. In characteristic two all permutation signs equal one,
so the permanent coincides with the determinant \cite{Berkowitz}.
The bound holds even for circuits of arbitrary depth that approximate
the desired coefficients and whose other output degrees are unrestricted.

\begin{theorem}[Unrestricted arithmetic circuits]\label{thm:main}
There is an absolute constant $c>0$ such that, for each characteristic
$\chi$ (including zero), there is an $N_0(\chi)$ with the following
property. Over every field $F$ of characteristic $\chi$ and at every
order $N\ge N_0(\chi)$, the degree-component border size satisfies
\[
  \bsz_F^{[N]}(\HC_N)\ \ge\
       \exp\!\left(c\,\frac{(\log N)^2}{(\log\log N)^2}\right).
\]
The same bound holds for $\Per_N$ if $\chi\ne2$.
\end{theorem}

\needspace{4\baselineskip}
Here $\HC_N$ and $\Per_N$ are forms of degree $N$ in the entries of
an $N\times N$ matrix $X$. A permutation $\pi$ contributes the
monomial $\prod_{i=1}^N X_{i,\pi(i)}$ to the permanent, including when
$\pi$ has fixed points. It contributes to $\HC_N$ exactly when $\pi$ has
a single cycle. Thus a tour is counted once, independently of its
starting vertex.

For a form $f$ of degree $D$, its \emph{degree-component border size}
$\bsz_F^{[D]}(f)$ is the least $s$ such that $f$ belongs to the Zariski
closure, over $\overline F$, of the degree-$D$ components of outputs
of circuits with at most $s$ gates. Each gate performs one binary
addition, subtraction or multiplication. Variables and scalar inputs
are free, and intermediate polynomials may have arbitrary degree.
We work with formal polynomials, including over finite fields. Since
this model permits limits as well as arbitrary output components in
other degrees, Theorem~\ref{thm:main} applies in particular to exact
circuits and to ordinary border complexity.

\subsection{Valiant's complete polynomials}

Valiant's class $\VP$ consists of polynomial families of polynomially
bounded degree and variable count computed by polynomial-size arithmetic
circuits. The class $\VNP$ allows, in addition, a sum over all assignments
to polynomially many auxiliary bits. Efficient verification of a cover,
followed by multiplication of its edge weights, puts both families in
$\VNP$. Hamiltonian cycle is complete for $\VNP$ over every field, and
the permanent is complete outside characteristic two. Every family in
$\VNP$ can therefore be obtained by substituting variables and constants
in the respective complete polynomial of polynomially larger order
\cite{Valiant,ValiantProjections,Burgisser}. This algebraic use of
projections is distinct from the counting complexity of evaluating a
permanent on a matrix of zeros and ones \cite{ValiantPermanent}.

In particular, $\VNP$ is not contained in border $\VP$ over any field
(see \cite[Conjecture~7.5]{BurgisserBorderConjecture}).
The companion paper on Valiant's hierarchy separates $\VP$ from $\VNP$
with a polynomial attached to a square grid \cite{oaH}. The bound in
Theorem~\ref{thm:main} gives a larger exponent for two classical
complete families and holds at every sufficiently large order. It is
already attained on sparse graphs. For Hamiltonian cycle one can retain
at most three incoming and three outgoing edges per vertex, and for the
permanent one can retain at most three entries per row and per column
(Corollary~\ref{cor:sparse-support}).

\subsection{Earlier lower bounds}\label{sec:earlier-bounds}
For exact computation of $\Per_N$ over $\mathbb C$, an earlier lower
bound of $\Omega(N^2\log\log N)$ counts binary arithmetic gates and
allows arbitrary depth and reuse of intermediate results
\cite[Chapter~5, Theorem~1.1]{TenAdvances}.
Under restrictions on the circuit, much larger bounds were known.
Jerrum and Snir proved exponential monotone circuit lower bounds for
both Hamiltonian cycle and the permanent
\cite[Sections~4.3 and~4.4]{jerrumsnir}.
Raz proved superpolynomial multilinear formula lower bounds for the
permanent and the determinant \cite{Raz}. Raz and Yehudayoff obtained
bounds exponential in a positive power of $N$ for both polynomials
at fixed depth in multilinear circuits \cite{RazYehudayoff}.
Depth alone can also force superpolynomial size, even for iterated
matrix multiplication \cite{LST,Forbes,BDS}.

A related question in geometric complexity theory \cite{MS} asks how
large an affine-linear matrix must be for its determinant to equal a
given polynomial, or to approach it coefficientwise. For the permanent,
Mignon and Ressayre proved a quadratic lower bound on the matrix order
in characteristic zero \cite{MignonRessayre}, and Cai, Chen and Li
obtained a quadratic bound in odd characteristic \cite{CaiChenLi}.
Landsberg, Manivel and Ressayre proved a quadratic bound allowing
limits over $\mathbb C$ \cite{LandsbergManivelRessayre}.
The corresponding cubic bound over $\mathbb C$ \cite{cubic-bdc} has
since been extended to each fixed odd characteristic \cite{Sheshadri}.
A determinant of order $m$ with affine-linear entries in $N^2$
variables has a division-free circuit of size polynomial in $m$ and
$N$ \cite{Berkowitz}. Consequently Theorem~\ref{thm:main} gives the
same superpolynomial lower bound on matrix order, up to the absolute
constant in the exponent, for Hamiltonian cycle in every characteristic
and for the permanent outside characteristic two.

\subsection{How the lower bound is proved}

Both proofs construct a hard polynomial as a weighted sum over
assignments accepted by a small Boolean verifier. Completeness
then transfers the bound to Hamiltonian cycle and the permanent. The
choice of witness turns on a property of the field characteristic that
comes from the geometry of tensors \cite{oaG}.

To see the geometry, imagine supplying a tensor by placing its indices
at sites and joining the sites by wires. A shared index with $q$
possible values has a cost proportional to $\log q$. Strassen's
asymptotic spectrum gives numerical lower bounds on the cost of
supplying a tensor in this way \cite{Strassen}. We call such a bound a
\emph{price}. For example, the logarithm of matrix rank across a fixed
cut charges the full $\log q$ for a wire across the cut and zero for a
wire on either side. In general the cost of the wire is its length
times $\log q$. The resulting distances give a class of metric spaces
for each characteristic (Section~\ref{sec:prices}).

A characteristic \emph{admits ports} when there are regular graphs of
bounded degree and arbitrarily large order with the following expansion
property. Place the vertices in any of its wire metric spaces. The
mean squared distance between two independent uniform vertices must be
at most a fixed multiple of the mean squared distance along a uniform
edge. The same constant must work for all the graphs and placements.
This is a nonlinear Poincar\'e inequality
\cite{MendelNaor,MendelNaorHadamard,EMN}. Either outcome gives a
circuit lower bound (Figure~\ref{fig:plan}). The proof requires no
classification of the characteristics.

\begin{figure}[!hb]
\centering
\begin{tikzpicture}[
  font=\small,align=center,
  branch/.style={draw=#1,fill=#1!4,rounded corners=2pt,line width=.55pt,
                 inner xsep=6pt,inner ysep=6pt,text width=5.25cm,minimum height=14mm},
  finish/.style={draw=ink,fill=ink!5,rounded corners=2pt,line width=.55pt,
                 inner sep=6pt,text width=9.0cm},
  arrow/.style={-{Stealth[length=1.8mm]},draw=ink,line width=.55pt}]
\node[branch=sea] (geom) at (-3.05,0)
  {\textbf{Without expansion}\\[2pt]
   One placement makes low-rank\\ terms cheap across every balanced cut.};
\node[branch=rust] (comp) at (3.05,0)
  {\textbf{With expansion}\\[2pt]
   Compression lets a program simulate\\ a proposed circuit for its own count.};
\node[branch=sea,below=5mm of geom] (binary)
  {An explicit binary tensor exceeds
   the price of these supplies.};
\node[branch=rust,below=5mm of comp] (signed)
  {A table of signed path counts forces a contradiction, even in the border model.};
\node[inner sep=0pt,outer sep=0pt,fit=(binary)(signed)] (witnesses) {};
\node[finish,below=6mm of witnesses] (both)
  {Each witness polynomial is a projection of Hamiltonian cycle,
   and of the permanent outside characteristic two.};
\draw[arrow] (geom)--(binary);
\draw[arrow] (comp)--(signed);
\draw[arrow] (binary.south)--(binary.south|-both.north);
\draw[arrow] (signed.south)--(signed.south|-both.north);
\end{tikzpicture}
\caption{The two witnesses for Theorem~\ref{thm:main}. Each verifier
gives a projection at polynomially larger order
(Lemma~\ref{gadget:verifier}). The counting argument also gives a
stretched-exponential bound in the characteristics where it applies
(Theorem~\ref{thm:stretch}).}
\label{fig:plan}
\end{figure}


\paragraph{When expansion fails.}
Consider a polynomial that uses one variable from each of $d$ rows in
every monomial. Its coefficients form a tensor with one index per row.
A small circuit expresses this tensor as a sum of terms, each of low
matrix rank across some balanced partition of the rows. The cut may
change from term to term. On four rows, take the product of wires on
$12$ and $34$, then add a second product on $13$ and $24$, using
disjoint local supports. Each term has rank one across its own cut.
Every balanced cut severs both wires of at least one term, so their
sum can have large rank across every cut.

Part~\ref{part:no} chooses \emph{one placement of the rows} at which
every low-rank term is cheap to supply. Start with all suppliers at the
center and the rows at unit distance. For each balanced cut, move two
suppliers towards their respective rows. A rank-$R$ term pays for one
wire of dimension $R$ between the suppliers, but saves on transport to
the rows. Every balanced cut saves order
$d/\sqrt{\log d}$ in transport per unit separation, to first order.
In Euclidean space some bisection of any centered set of unit vectors
gains only $O(\sqrt d)$, by averaging over random bisections. Failure of
expansion yields directions with the larger gain through an entropy
inequality for random walks. A product construction combines these
directions for all cuts (Sections~\ref{sec:poincare}--\ref{sec:lamps}).

At those same sites we need a tensor that costs nearly as much as
sending all its indices from the center. Exterior multiplication gives
tensors of minors whose prices can be estimated through dimensions of
intersections of subspaces. We construct one explicit tensor
with coefficients in $\{0,1\}$ from which a suitable tensor of minors
can be obtained at every such placement
(Sections~\ref{sec:exterior}--\ref{sec:packing}). Its price exceeds
the price of a sum of cheap supplies, even after accounting for the wire
that tells every row which summand to use. The comparison also works on
every grouped tensor power, keeping the degree at $d$ while the row
length grows. For fixed sufficiently large even $d$, this yields circuit size
$n^{\Omega(d/\sqrt{\log d})}$ in a variable budget $n$
(Corollary~\ref{cor:variable-budget}).
Letting $d$ grow gives the exponent of Theorem~\ref{thm:main}.

\paragraph{When expansion holds.}
In Part~\ref{part:port} the same geometry yields identities that
compress a tensor network while preserving its value \cite{oaG}.
Here a coefficient of the witness is the signed number of computation
paths of a program on a specified input. A sufficiently small network
for its coefficient table would let a program use compression to
simulate its own count within its allotted time and add one. The
contradiction depends on which work the program must do itself and which
it can delegate to the branches it counts. Sections~\ref{sec:compr}--%
\ref{sec:batches} establish this distinction and give a lower bound
even for networks computing batches of independent copies of the table
over a finite field (Theorem~\ref{thm:batches}). This use of an
algorithm to prove hardness has precedents in circuit complexity
\cite{ki,williams} and in tensor-network algorithms for satisfiability
\cite{falseeth,Anthropic}.

A border circuit poses a further problem. A coefficient that has a
finite boundary value can be a sum of terms with poles. On an algebraic
curve through the boundary point we use residues and the Cartier
operator to recover the coefficient as an eigenvalue of an explicitly
accessible matrix. The program can now perform its counting
contradiction one matrix entry at a time. For characteristic zero,
specialization to a suitable prime preserves the finite data needed by
the simulation (Sections~\ref{sec:parse}--\ref{sec:border}).
In every characteristic $\chi$ admitting ports, this argument gives
the stronger bound $2^{N^{\varepsilon_\chi}}$, for some
$\varepsilon_\chi>0$, on $\HC_N$ and, when $\chi\ne2$, on $\Per_N$
(Theorem~\ref{thm:stretch}).
